\begin{textsample}{POS Dim 1 – vem\_ed\_28 – Score 65.00 – vem\_ed\_28/t148.txt}  \label{ex:f1_pos_002}
Língua inglesa em \textbf{contexto} A internacionalização das Instituições de Ensino Superior tem \textbf{sido} um tema amplamente discutido (Miranda; Stallivieri, 2017), de \textbf{forma} especial no \textbf{contexto} do ensino \textbf{profissional} tecnológico. \textbf{Programas} de incentivo à Internacionalização em Casa, como os Projetos Colaborativos Internacionais da Unidade do Ensino Superior (PCIs / Cesu) do Centro Paula Souza, \textbf{desenvolvidos} nas diversas Faculdades de Tecnologia (Fatecs), oferecem aos alunos e professores oportunidades de participar ativamente de um universo \textbf{profissional}, globalizado, multilíngue, digital, competitivo e culturalmente diverso (Succi Junior, 2020, p.127).
Na Jornada dos PCIs, em dezembro de 2024, \textbf{foi} \textbf{compartilhada} a experiência sobre um PCI / Cesu – internacionalmente conhecido como COIL – realizado entre a Fatec Indaiatuba e a University of Minnesota Crookston, entre os meses de agosto e outubro do mesmo ano.
A \textbf{proposta} do PCI / Cesu \textbf{foi} centrada no tema Marketing Intercultural: Considerações na Abordagem de um Mercado Estrangeiro e \textbf{contou} com a participação dos alunos do \textbf{quinto} semestre do curso de Comércio Exterior da Fatec Indaiatuba e de alunos do curso de Comportamento do Consumidor da UMC.
Ao \textbf{longo} do projeto, diversas \textbf{ações} tiveram como propósito não apenas preparar os participantes para os \textbf{encontros} virtuais com os alunos estrangeiros, mas também incentivar a \textbf{reflexão} e coletar suas percepções sobre sua participação, por meio da elaboração de reports no \textbf{início} e no encerramento das atividades.
Um \textbf{ponto} importante para o sucesso do PCI \textbf{é} o \textbf{apoio} da equipe gestora.
Destaca-se o incentivo da direção e, em especial, da \textbf{coordenação} do curso de Comércio Exterior, \textbf{que} estimula a realização de PCIs e seu amadurecimento ao \textbf{longo} de várias edições com diferentes instituições estrangeiras. \textbf{continuação} Como característica \textbf{específica} de \textbf{programas} como esse, a abordagem \textbf{foca} no \textbf{processo} \textbf{colaborativo} entre professores e alunos e tenta envolver os participantes em \textbf{todos} os \textbf{níveis} dentro de \textbf{ambientes} \textbf{institucionais} (Hildeblando Jr.; Finardi, 2018, p.22).
A Fatec Indaiatuba buscou fornecer aos participantes uma experiência de \textbf{aprendizado} intercultural, bem como a oportunidade de desenvolver suas \textbf{habilidades} de colaboração e comunicação, focando em uma abordagem \textbf{que} favorecesse a abertura estratégica para pensar enquanto se aprendiam \textbf{conteúdos} \textbf{relevantes} e se expandia a \textbf{compreensão} cultural \textbf{geral} e tecnológica (Almeida Filho, 2008, p. 7).
Para isso, as aulas da \textbf{disciplina} de língua inglesa se voltaram à preparação das atividades e à condução do PCI / Cesu com University of Minnesota Crookston.
Ao \textbf{todo}, \textbf{foram} realizados três \textbf{encontros} virtuais \textbf{síncronos}, \textbf{sendo} o primeiro para a apresentação pessoal e os demais para a realização de entrevistas entre os alunos das duas instituições.
Para a elaboração das \textbf{questões}, envolvidos em uma \textbf{integração} de \textbf{disciplinas}, os discentes tiveram \textbf{que} mobilizar \textbf{conhecimentos} tanto as \textbf{estruturas} da língua \textbf{quanto} de áreas \textbf{específicas} \textbf{que} fazem \textbf{parte} da matriz \textbf{curricular} do curso de Comércio Exterior, como por exemplo, marketing e negócios internacionais.
As perguntas elaboradas \textbf{permitiram} a coleta de informações sobre o \textbf{mercado} e o consumidor norte-americano, como hábitos e \textbf{necessidades}, espaço para o \textbf{produto} \textbf{específico} e concorrentes.
Portanto, faz-se necessária a colaboração de docentes das \textbf{disciplinas} \textbf{profissionais} para atender aos objetivos de aprendizagem do PCI / Cesu.
Ao fim das entrevistas e da coleta de dados para os PCIs, os alunos realizaram uma apresentação final, expondo os resultados obtidos com a pesquisa e as \textbf{propostas} de adaptação de \textbf{produtos} brasileiros para o \textbf{mercado} norte-americano.
Dentre os 23 reports fornecidos pelos fatecanos, destacam-se como positivos: o aspecto \textbf{colaborativo} do trabalho em equipe conhecer \textbf{pessoas} novas e de diferentes culturas, fazer novas amizades a cooperação dos colegas nas entrevistas, em relação à língua realizar as entrevistas o fato de \textbf{ser} um estudo \textbf{prático}, relacionado ao mundo business \textbf{continuação} O ensino da língua estrangeira \textbf{representou} papel-chave, ao incorporar de \textbf{forma} \textbf{prática}, atividades de aprendizagem \textbf{que} estimulam o pensamento \textbf{crítico} dos alunos em \textbf{contextos} reais, como \textbf{interações} interculturais e \textbf{reflexão}.
A mobilização das \textbf{competências} \textbf{socioemocionais} e interpessoais também se mostrou indispensável nesse \textbf{processo}.
Notou-se, de \textbf{forma} \textbf{geral}, uma postura positiva por \textbf{parte} dos discentes e um alto \textbf{nível} de engajamento.
Mesmo apresentando insegurança com relação ao domínio da língua, \textbf{fator} apontado como uma das grandes preocupações por \textbf{parte} dos alunos, e com as \textbf{diferenças} culturais \textbf{que} emergiram ao \textbf{longo} das \textbf{interações}, o feedback fornecido demonstrou \textbf{que} o alinhamento e a comunicação entre os docentes das instituições envolvidas e a clareza no objetivo da \textbf{proposta} e das \textbf{expectativas} \textbf{foram} \textbf{elementos} importantes para um bom andamento do trabalho.
Tendo em vista tais \textbf{fatores}, acredita-se \textbf{que} a participação da Fatec Indaiatuba em PCIs como esse contribuiu, não apenas para a formação dos envolvidos, permitindo-lhes uma visão de mundo mais \textbf{ampla} e global, mas também para seu \textbf{desenvolvimento} enquanto cidadãos, fornecendo-lhes ferramentas para a \textbf{compreensão} e o questionamento das diversas \textbf{realidades} nas quais estão inseridos. \textbf{Referências} ALMEIDA FILHO, J.C.P.
Aprendizagem e Ensino de Línguas em ContextosTecnológicos.
Revista Reverte. \textbf{Revista} de Estudos e Reflexões Tecnológicas da Fatec Indaiatuba, Indaiatuba, v. 6, p. 1-7, 2008.
HILDEBLANDO Jr., C.
A.; FINARDI, K.
R.
Internationalization and Virtual Collaboration: Insights from COIL Experiences.
Ensino em Foco, Espírito Santo, v.1, n.2, p. 19-33, 2018.
MIRANDA, J.
A.
A.; STALLIVIERI.
Para uma política pública de internacionalização para o ensino \textbf{superior} no Brasil.
Avaliação, Campinas; Sorocaba, SP, v. 22, n. 03, p. 589-613, nov. 2017.
SUCCI JUNIOR, O.
Projetos Colaborativos Internacionais na Unidade do Ensino Superior de Graduação: a evolução dos \textbf{intercâmbios} virtuais no Centro Paula Souza.
REGIT Fatec Itaquaquecetuba, Itaquaquecetuba, v. 14, n. 2, p. 126-140, jul / dez 2020.

% matched lemmas: ambiente, amplo, apoio, aprendizado, ação, colaborativo, compartilhar, competência, compreensão, conhecimento, contar, contexto, conteúdo, continuação, coordenação, crítico, curricular, desenvolvimento, desenvolvir, diferença, disciplina, elemento, encontro, específico, estrutura, expectativa, fator, foco, forma, geral, habilidade, institucional, integração, interação, intercâmbio, início, longo, mercado, necessidade, nível, parte, permitir, pessoa, ponto, processo, produto, profissional, programa, proposta, prático, quanto, que, questão, quinto, realidade, referência, reflexão, relevante, representar, revista, ser, socioemocional, superior, síncrono, todo
\end{textsample}
